%! Problems with Sampling — Unit 3, Lesson 3.3 — Lesson Plan
% GRADUAL-RELEASE plan (warm-up / guided notes and practice / debrief / close and start the
% homework) at the COMPACT budget: the whole lesson is finished in one 55-minute period.
% THERE IS NO GROUP ACTIVITY IN THIS COURSE and NO INDEPENDENT PRACTICE SET in the notes — the
% homework is the individual practice, started in class. Teacher-facing. Every box is `skillbox`
% (breakable) with a \boxguard ahead of it; no `fixedskillbox`, no tier boxes, no exit ticket.
% CED Topic 3.4. Lesson 3.2 taught the random sampling methods; this lesson names what goes
% wrong anyway. Lesson 3.4 turns to experiments.
\documentclass[10pt]{article}
\usepackage{apstats-article}
\usepackage{apstats-boxes}

\newcommand{\UnitNumberName}{Unit 3: Collecting Data \quad}
\newcommand{\LessonNumberName}{Lesson 3.3: Problems with Sampling}

\begin{document}
\begin{center}
    {\Large\bfseries \CourseName \\
    {\small\bfseries \UnitNumberName \LessonNumberName}}
\end{center}

\begin{tcolorbox}[colback=sky, colframe=navy, boxrule=0.9pt, arc=2mm,
    left=2mm, right=2mm, top=1.5mm, bottom=1.5mm]
    \textbf{Primary Objective:} Students will identify voluntary response, convenience,
    undercoverage, nonresponse, and response bias in a described sampling method; for each,
    explain in context who is over- or under-represented and whether the estimate is likely
    too high or too low; and reject the claim that a larger sample repairs a biased method.
    \\[2pt]
    \textbf{CED alignment:} Topic 3.4 (Potential Problems with Sampling) \quad$\cdot$\quad Big
    Idea: DAT (Data-Based Inferences) \quad$\cdot$\quad Skill 1.C \quad$\cdot$\quad LO DAT-2.E
    \quad$\cdot$\quad EK DAT-2.E.1 (bias: responses systematically favored), DAT-2.E.2
    (voluntary response), DAT-2.E.3 (undercoverage), DAT-2.E.4 (nonresponse), DAT-2.E.5
    (response bias: question wording, self-report), DAT-2.E.6 (convenience and voluntary
    response samples do not use chance)
    \\[2pt]
    \textbf{Lesson model:} \emph{Gradual release --- I do, we do, you do --- all of it inside
    the guided notes.} There is no group activity. The notes name the vocabulary as they teach
    it and deliver the instruction row by row in one \emph{Main Ideas / Notes} table --- each
    idea a definition to complete and a grid of short problems, the first worked by you and the
    rest by the class; the last row, Guided Practice, is worked together with students holding
    the pen; \textbf{the homework is the individual practice}, and students start it in class
    while you circulate. The whole-class debrief is \emph{spoken}.
    \\[2pt]
    \textbf{Compact budget:} the notes are three pages, 14 problems. Every problem asks the same
    three things --- name the bias, say who it favors, give the direction --- so the habit is
    built by repetition. SRS, stratified, and cluster samples are Lesson 3.2's: assume them, do
    not reteach them.
\end{tcolorbox}

\renewcommand{\arraystretch}{1.4}
\boxguard
\begin{skillbox}[Learning Targets \& Key Understandings]{goldbox}
\begin{tabularx}{\linewidth}{@{}X|X@{}}
\textbf{Learning targets (student language)} & \textbf{Key understandings (the ``why'')} \\[2pt]
\hline
\vspace{-6pt}
\begin{itemize}[leftmargin=1.1em, nosep, after=\vspace{2pt}]
  \item I can identify voluntary response bias and the bias in a convenience sample, and say
        which way each pushes the estimate.
  \item I can tell undercoverage bias from nonresponse bias, even when the sample was chosen
        at random.
  \item I can explain how a question's wording or a self-reported answer creates response
        bias.
  \item I can explain why a larger sample reduces variability but does not fix bias.
\end{itemize}
&
\vspace{-6pt}
\begin{itemize}[leftmargin=1.1em, nosep, after=\vspace{2pt}]
  \item Bias occurs when certain responses are systematically favored over others
        (DAT-2.E.1). Volunteers and convenience samples do not use chance, so they invite it
        (DAT-2.E.2, DAT-2.E.6).
  \item A random sample can still be biased: the list can leave people out (undercoverage,
        DAT-2.E.3), the chosen can fail to answer (nonresponse, DAT-2.E.4), and the question
        or the asker can push the answers (response bias, DAT-2.E.5).
  \item \textbf{The AP habit:} naming the bias is half an answer. A full answer says who is
        over- or under-represented and whether the estimate is too high or too low, in
        context. The direction comes from \emph{who} answers, not from the bias's name ---
        voluntary response can push an estimate down (problem~3).
  \item \textbf{Target misconception:} ``a huge sample fixes bias.'' A larger sample reduces
        variability, not bias. 48{,}000 forum volunteers said 81\%; 600 players at random
        said 47\%. The 600 win.
\end{itemize}
\end{tabularx}
\end{skillbox}

\boxguard
\begin{skillbox}[{Vocabulary, Concepts, \& Theorems --- taught directly in the guided notes}]{greenbox}
\begin{minipage}[t]{0.48\linewidth}
    \begin{itemize}
        \item \textbf{Voluntary response bias} --- people choose to respond; those with strong
              opinions are over-represented (DAT-2.E.2)
        \item \textbf{Convenience sample} --- whoever is easiest to reach; like a voluntary
              response sample, it does not use chance (DAT-2.E.6)
        \item \textbf{Undercoverage bias} --- part of the population has a reduced or no chance
              of being chosen (DAT-2.E.3)
    \end{itemize}
\end{minipage}
\hfill
\begin{minipage}[t]{0.48\linewidth}
    \begin{itemize}
        \item \textbf{Nonresponse bias} --- chosen individuals can't be reached or refuse, and
              may differ from those who answer (DAT-2.E.4)
        \item \textbf{Response bias} --- the question or the way it is asked pushes answers:
              leading or confusing wording, self-reported answers (DAT-2.E.5)
        \item \textbf{Size vs.\ bias} --- a larger sample reduces variability, not bias
        \item \textbf{Carried forward} --- SRS, stratified, and cluster samples (Lesson 3.2)
    \end{itemize}
\end{minipage}
\end{skillbox}

% A tabularx never splits, so guard generously ahead of it.
\boxguard[30]
\begin{skillbox}[Lesson at a Glance --- \MeetingLength]{sky}
\renewcommand{\arraystretch}{1.25}
\begin{tabularx}{\linewidth}{@{}l c X X@{}}
\textbf{Phase} & \textbf{Min} & \textbf{Students} & \textbf{Teacher} \\
\hline
Warm-Up & 5 & Name two methods from 3.2, Maya's lunch-table survey, and two wordings of one
question & Leave ``who is left out?'' (item~2) and ``(B) gets more yes'' (item~3) on the
board \\
Guided Notes \& Practice & 34 & Fill the vocabulary box and the notes table: problems 1--10 row
by row, then \emph{The Work-Hours Survey} (11--14) together, pen in their hand & \textbf{I do}
each row's definition lines and first problem (1, 4, 7), \textbf{we do} the rest (24)
$\to$ \textbf{we do} Guided Practice (10) \\
Debrief & 8 & Pens down, papers up --- answer aloud, nothing to fill in & Put 81\% (48{,}000
volunteers) beside 47\% (600 at random); open on the almost-right problem~9 \\
Close \& Assign & 8 & \textbf{Start the homework in class}, alone and silent & Circulate on
Part~B item~2 for your formative read; preview Lesson 3.4 \\
\end{tabularx}
\end{skillbox}

\boxguard[20]
\begin{skillbox}[Warm-Up --- Activate Prior Knowledge (5 min)]{sky}
\begin{minipage}[t]{0.48\linewidth}
    \textbf{The items and what each seeds}
    \begin{itemize}
        \item \textbf{1.} Name two methods: a computer picks 60 of 1{,}200 (SRS); 15 at random
              from each grade (stratified). \textbf{Spirals} to Lesson 3.2 and sets up the
              point that a random method is the baseline today's biases spoil.
        \item \textbf{2.} Maya asks the 12 friends at her lunch table. \textbf{Seeds} rows 1
              and 2: the room's ``she just asked her friends'' becomes a \emph{convenience
              sample}, and ``who is left out'' becomes \emph{undercoverage}.
        \item \textbf{3.} Two wordings of a phone-ban question. \textbf{Seeds} row~3: same
              students, different wording, different answer --- \emph{response bias}.
    \end{itemize}
\end{minipage}
\hfill
\begin{minipage}[t]{0.48\linewidth}
    \textbf{Running it}
    \begin{itemize}
        \item \textbf{Item~1 is a 60-second spiral}; take answers and move on. If the room
              cannot name stratified, note it for Lesson 3.2's reteach, not today.
        \item \textbf{Debrief item~2 aloud} and write the room's answer to ``who is left out?''
              on the board in their words. Rows 1 and 2 name those ideas.
        \item \textbf{Take a hand count on item~3} and leave ``(B) gets more yes --- why?'' on
              the board. Row~3 points back at it.
    \end{itemize}
\end{minipage}
\end{skillbox}

\boxguard
\begin{skillbox}[Guided Notes \& Practice --- I do, then we do (34 min)]{sky}
\textbf{This is the whole lesson.} There is no group activity, and there is no independent
practice set on this page: \textbf{the homework is the individual practice}, started in class.
The notes are one \emph{Main Ideas / Notes} table --- four rows, 14 short problems.
\textbf{In every row you fill the definition lines and work the first problem of each grid; the
class works the rest, pen in their hand.} The vocabulary box is filled \emph{as you go}.
\begin{multicols}{2}
\textbf{Volunteers (7 min).} Fill the definitions (choose / strong / easiest / chance) and the
three-step \emph{complete answer} --- it is the procedure for every problem today. Work
\textbf{problem~1} yourself, out loud, all three steps: voluntary response, frustrated drivers
vote, too high. Problem~2 is theirs. \textbf{Problem~3 is the trap}: students who have just
heard ``too high'' will say it again. Press: \emph{``Who bothers to fill out a comment card?''}
Unhappy diners --- so the satisfied share is too \emph{low}. The direction comes from who
answers, not from the name of the bias.

\textbf{Coverage and Response (7 min).} Fill the definitions and read \emph{the test} aloud:
never on the list $\to$ undercoverage; picked but silent $\to$ nonresponse. Work
\textbf{problem~4} (the landline book misses cell-only residents). Problems 5 and 6 are
theirs; 6 is the quiet one --- the principal picked \emph{at random}, but from a list that
leaves out students at work. Say it: \emph{random is not a vaccine}.

\textbf{Response Bias (10 min) --- the crux.} Point back at warm-up item~3 on the board, fill
the definition, and work \textbf{problem~7} (the turf field's safety reason). Problem~8 is
theirs. Then read the game-forum context and \textbf{take a hand count on problem~9
\emph{before} anyone explains} --- that is the crux. Problem~10 resolves it: chance picked the
600. Fill the callout --- \emph{a larger sample reduces variability, not bias} --- and make the
room say it back.

\columnbreak
\textbf{We do (about 10 min): The Work-Hours Survey.} A fresh context, problems 11--14, worked
entirely together. One plan, three steps, three different biases --- all pushing the same way.
Students hold the pen; you hold the questions: \emph{``Could an off-campus student ever be
picked?''} \quad \emph{``Who skips a meeting with their advisor?''} \quad \emph{``Would you
tell your advisor you work 30 hours?''} \textbf{Problem~14 is the crux transferred}: 6{,}000
emailed replies against the rest. It is the one to protect when time is short.

\textbf{The pen must actually change hands.} Put problem~11 on them cold, take a hand count on
problem~14 before anyone speaks, and make the room defend the answer. Circulate and demand the
direction:
\begin{itemize}[leftmargin=1.1em, nosep]
  \item \textbf{Problem~11:} ``Too high or too low --- and who is missing?''
  \item \textbf{Problem~12:} ``Were these 500 chosen at random? Then why is it still biased?''
  \item \textbf{Problem~14:} ``If all 30{,}000 replied, would we be done? And if 6{,}000
        chose to?''
\end{itemize}
While circulating, pick the papers to put up in the debrief --- one \emph{almost}-right
problem~9 (``disagree, because 600 is enough'') and one that names a bias with no direction.
\textbf{Your formative read comes twice}: here, and again in the last eight minutes as students
start the homework.
\end{multicols}
\end{skillbox}

\boxguard
\begin{skillbox}[Debrief --- whole class (8 min)]{sky}
Pens down, papers up. \textbf{This debrief is spoken} --- nothing on the page to fill in.
\begin{multicols}{2}
\begin{enumerate}[leftmargin=1.3em, nosep, itemsep=3pt]
  \item \textbf{Open on the almost-right problem~9}: ``disagree, because 600 is enough.'' Right
        verdict, wrong reason --- it argues about size, the one thing that does not matter.
        Correct it in public: the forum is still volunteers; size shrinks variability, it
        does not move the center.
  \item Put a bias-with-no-direction answer up (``it's undercoverage''). Ask the room to finish
        it: who is missing, and too high or too low? That is the AP standard.
  \item Take problem~3 back: why did the comment cards lean \emph{low} when the traffic poll
        leaned \emph{high}? Same bias, different people who care.
\end{enumerate}
\columnbreak
\textbf{The sentence to leave on every desk:} \emph{Name the bias, say who it favors, and say
which way it pushes. A larger sample reduces variability, not bias.} Make a student say it,
not you.

\textbf{If the debrief runs short}, cut item~3 --- item~1 is the only one that cannot be cut.
\textbf{Do not let the debrief eat the last eight minutes} --- the homework start is your
second formative read.
\end{multicols}
\end{skillbox}

\boxguard
\begin{skillbox}[AP Practice --- extra credit, scored 1 to 5]{goldbox}
Two multiple-choice items --- \textbf{MC~1} names undercoverage \emph{and} its direction in a
transit survey drawn from vehicle owners (B; C has the right name with the wrong direction, E
is ``random, so no bias''), and \textbf{MC~2} is a leading speed-limit question (B: response
bias, too high; A is ``random, so unbiased'' and E is the crux --- ``2{,}000 would remove the
bias'') --- then one free-response set on a county sleep survey: \textbf{(a)} a source of bias
other than wording (nonresponse: 95 of 400, busy students sleep less and reply less, so the mean
is too high), \textbf{(b)} the wording's push (``doctors say 8 to 10 hours'' inflates reports),
and \textbf{(c)}, \textbf{the crux}: emailing all 9{,}000 does not fix either bias.

\textbf{Extra credit, scored 1 to 5, turned in with the homework.} Go over it at
the start of Lesson 3.4. Hold the AP standard out loud: \emph{a bias with no direction has not
earned the point} --- part (a) needs who is missing, why they differ, and too high or too low.
\end{skillbox}

\boxguard
\begin{skillbox}[Homework --- scored 1 to 5, due the first class after two study halls]{goldbox}
Five items in a third context: the state DOT estimating how many licensed teen drivers use a
phone while driving. \textbf{Part~A}: \textbf{A1} a social-media poll (voluntary response; the
safety-minded answer, so 22\% is too low); \textbf{A2} a trooper phoning a random sample
(response bias in a random sample, too low); \textbf{A3} 390 of 1{,}000 never pick up
(nonresponse, and how they differ). \textbf{Part~B}: \textbf{B1} an AP-style multiple choice ---
which change reduces A2's bias (B: an anonymous form; A, a bigger sample, is the distractor) ---
with a required justification; \textbf{B2} the crux restated: ``3{,}000 responses, three times
as many, so Plan~1 is more trustworthy.''

\textbf{Start it in the last eight minutes of the period, in class and alone.}
\textbf{Scoring:} 10 points --- A1, A2, A3, B2 at 2 each; B1 1 for the letter and 1 for the
justification. \textbf{B2 is the one to read closely}: a student who sides with the reporter has
kept the target misconception. The spiralbox previews Lesson 3.4.

\textbf{Paper homework is the default for this course} --- assign this page unless you have
decided otherwise. \textbf{DeltaMath override (occasional):} on the days you assign DeltaMath
instead, it replaces this page, and this page stays in the packet as extra practice. Never
assign both.
\end{skillbox}

\boxguard
\begin{skillbox}[Watch For (while circulating)]{redbox}
\begin{itemize}[leftmargin=1.2em, nosep]
  \item \textbf{``A huge sample fixes bias''} (notes~9 and 14, HW~B2, AP~MC2 (E) and (c) --- the
        target misconception). Probe: \emph{``If 48{,}000 people volunteered, who would they
        be? The same kind of people, only more of them.''}
  \item \textbf{A bias named with no direction} (every problem; HW~A1--A3, AP~(a)). Probe:
        \emph{``Too high or too low --- and who pushed it there?''}
  \item \textbf{``Voluntary response always makes it too high''} (notes~3, HW~A1). Probe:
        \emph{``Who cares enough to answer? Which way do \emph{they} lean?''}
  \item \textbf{``It was random, so it can't be biased''} (notes~6 and 12, HW~A2, AP~MC1 (E)
        and MC2 (A)). Probe: \emph{``Random from which list? And did everyone chosen answer,
        honestly?''}
  \item \textbf{Mixing up undercoverage and nonresponse} (notes~4--6, 11--12). Probe:
        \emph{``Could this person ever have been picked?''}
\end{itemize}
Cold-call prompts: \emph{``Give me a random sample that is still biased.''} \quad
\emph{``Would you rather have 600 at random or 48{,}000 volunteers? Why?''}
\end{skillbox}

\boxguard
\begin{skillbox}[Close \& Assign (8 min)]{goldbox}
\textbf{Students start the homework here, in class, alone and silent --- this is the individual
practice, and the first minutes of it are your best formative read.} Hand the launch first ---
five items on paper, scored, due the first class after two study halls --- and point out that
the AP Practice page is extra credit, turned in with the homework. Circulate on \textbf{Part~B
item~2}, then glance at A1's direction. Sort what you see into three piles as you walk ---
disagreed and said 3{,}000 volunteers are still volunteers (ready); disagreed with no reason,
or ``because 1{,}000 is enough'' (ready, needs the language); agreed with the reporter (the
misconception survived) --- and open Lesson 3.4 by returning to the 81\%/47\% board if that
third pile ran past a third of the room. Then close on what changed: \emph{``You came in
naming methods. You leave able to say how a method goes wrong --- and which way.''}
\textbf{Preview:} Lesson 3.4, \emph{Designing Experiments} --- treatments, random assignment,
control, and why only an experiment can show cause.
\end{skillbox}

% ── Teacher notes — one per component, in packet order (four: there is no activity) ──
\begin{teachernote}[Warm-Up]
Five minutes total. Item~1 is a 60-second spiral to Lesson 3.2 --- do not reteach the methods
here. Items 2 and 3 are the seeds: write the room's answers to ``who is left out?'' and ``why
does (B) get more yes?'' on the board in their words; rows 1--3 of the notes name both ideas
and point back at them.
\end{teachernote}

\begin{teachernote}[Guided Notes \& Practice]
\textbf{Pace it 24 / 10, then 8 for the debrief, then 8 for the homework start.} Four rows, 14
problems. \textbf{You work 1, 4, and 7}; the class works the rest. Volunteers and Coverage must
move --- seven minutes each --- because the time belongs to the Response Bias row, which
carries the crux. \textbf{Problem~3 is the trap} (voluntary response pushing an estimate
\emph{down}); \textbf{problem~9 is the crux} (48{,}000 forum volunteers ``must be close to the
truth''). Take a hand count on 9 before any explanation; problem~10 and the callout are the
resolution. \textbf{The Work-Hours Survey is where the pen changes hands}; problem~14 is the
crux transferred and the one to protect. If you only get through 11, 12, and 14, that is the
right trade. \textbf{Your second formative read is the homework start}: Part~B item~2, three
piles. \textbf{Debrief (8 min), spoken}, opening on the almost-right problem~9.
\end{teachernote}

\begin{teachernote}[AP Practice]
\textbf{This page is extra credit, scored 1 to 5} --- the graded practice is the homework
behind it. \textbf{MC~1 is the one worth reading closely}: (C) has the right bias with the
wrong direction, which is exactly the AP habit this lesson builds; if many choose it, open
Lesson 3.4 with problem~3's comment cards. On the free response, part (a) must name the bias,
say how the non-responders differ, \emph{and} give a direction; part (c) must say why a bigger
sample leaves the bias in place, not just ``it's still biased.''
\end{teachernote}

\begin{teachernote}[Homework]
\textbf{Scored 1 to 5 (10 points), due the first class after two study halls.} The eight minutes
at the end of the period are a \emph{start}: put students on B2 and A1 first while you can
still catch a wrong start. \textbf{B2 is the item to read first when scoring} --- a student who
sides with the reporter has kept the target misconception. On A1, accept a well-argued ``too
high,'' but a direction with no reason about who answers a DOT poll earns 1 of 2. On A2, the
bias must be response bias (the sample is random) with the direction too low. On A3, the way
the 390 differ must bear on phone use while driving. On B1, (B) with no justification earns
1 of 2. \textbf{Paper is the default.} On the occasions you assign DeltaMath in its place, tell
students this page is now optional practice. Do not assign both.
\end{teachernote}
\end{document}
